% Part: normal-modal-logic
% Chapter: axioms-systems
% Section: introduction

\documentclass[../../../include/open-logic-section]{subfiles}

\begin{document}

\olfileid[pt]{nml}{axs}{int}

\olsection{Introdução}

Temos uma semântica para a linguagem modal básica em termos de modelos
modais, e uma noção de !!a{formula} ser válida---verdadeira em todos os mundos
de todos os modelos---ou válida com respeito a alguma classe de modelos ou
frames---verdadeira em todos os mundos de todos os modelos da classe, ou
baseados no frame. A lógica normalmente conecta essas caracterizações
semânticas da validade a uma noção teórica da prova de !!{derivability}. O
objetivo é definir uma noção de !!{derivability} em algum sistema tal que
!!a{formula} seja !!{derivable} se e somente se for válida.

Os sistemas de !!{derivation} mais simples e historicamente mais antigos são
os chamados sistemas de !!{derivation} de tipo Hilbert ou axiomáticos. Os
sistemas de !!{derivation} de tipo Hilbert para muitas lógicas modais são
relativamente fáceis de construir: são simples como objetos de estudo
metateórico (p.ex., para demonstrar correção e completude). Contudo, é
muito mais difícil usá-los para demonstrar !!{formula}s do que, por exemplo,
usar sistemas de dedução natural.

Nos sistemas de !!{derivation} de tipo Hilbert, !!a{derivation} de
!!a{formula} é uma sequência de !!{formula}s que parte de certos axiomas e,
por meio de algumas regras de inferência, chega à !!{formula} em questão. Como
queremos que o sistema de !!{derivation} corresponda à semântica, precisamos
garantir que o conjunto de fórmulas !!{derivable} seja verdadeiro em todos os
modelos (ou verdadeiro em todos os modelos nos quais todos os axiomas são
verdadeiros). Primeiro isolaremos algumas propriedades das lógicas modais
necessárias para que isso funcione: as lógicas modais «normais». Para lógicas
modais normais, só é preciso supor duas regras de inferência: modus ponens e
necessitação. Como axiomas, tomamos todas as instâncias de substituição de
tautologias e, dependendo da lógica modal de que tratamos, vários axiomas
modais. Mesmo que estejamos interessados apenas na classe de todos os modelos,
também devemos contar como axiomas todas as instâncias de substituição
de~$\Ax{K}$\iftag{notprvDiamond}{}{ e de~$\Ax{Dual}$}. Isso, por si só, gera
a lógica modal normal mínima~$\Log K$.

\begin{defn}
A regra de \emph{modus ponens} é o esquema de inferência
\begin{prooftree}
\AxiomC{$!A$}
\AxiomC{$!A \lif !B$}
\RightLabel{\MP}
\BinaryInfC{$!B$}
\end{prooftree}
Dizemos que !!a{formula}~$!B$ \emph{segue de}~!!{formula}s $!A$, $!C$ por
modus ponens se e somente se $!C \ident !A \lif !B$.
\end{defn}

\begin{defn}
A regra de \emph{necessitação} é o esquema de inferência
\begin{prooftree}
\AxiomC{$!A$}
\RightLabel{\Nec}
\UnaryInfC{$\Box !A$}
\end{prooftree}
Dizemos que a !!{formula}~$!B$ segue das !!{formula}s $!A$ por necessitação se
e somente se $!B \ident \Box !A$.
\end{defn}

\begin{defn}
Uma \emph{!!{derivation}} a partir de um conjunto de axiomas~$\Sigma$ é uma
sequência de !!{formula}s $!B_1$, $!B_2$, \dots, $!B_n$, na qual cada $!B_i$
é
\begin{enumerate}
\item uma instância de substituição de uma tautologia, ou
\item uma instância de substituição de !!a{formula} em~$\Sigma$, ou
\item segue de duas !!{formula}s $!B_j$, $!B_k$, com $j$, $k < i$, por modus
  ponens, ou
\item segue de !!a{formula}~$!B_j$, com $j < i$, por necessitação.
\end{enumerate}
Se existe tal !!a{derivation} com $!B_n \ident !A$, dizemos que $!A$ é
\emph{!!{derivable} de $\Sigma$}, em símbolos, $\Sigma \Proves !A$.
\end{defn}

Com essa definição, veremos que o conjunto de fórmulas !!{derivable} forma uma
lógica modal normal, e que toda expressão !!{derivable} que seja uma
!!{formula} é verdadeira em todo
modelo no qual cada axioma é verdadeiro. Essa propriedade das
!!{derivation}s é chamada de \emph{correção}. A recíproca, a
\emph{completude}, é mais difícil de demonstrar.

\end{document}
